\documentclass[11pt]{amsart}

\usepackage[margin=1.15in]{geometry}
\usepackage{amsmath,amssymb,amsthm,mathtools}
\usepackage{microtype}
\usepackage{hyperref}
\usepackage[nameinlink,capitalize]{cleveref}

\hypersetup{
  colorlinks=true,
  linkcolor=blue,
  citecolor=blue,
  urlcolor=blue
}

\theoremstyle{plain}
\newtheorem{theorem}{Theorem}[section]
\newtheorem{proposition}[theorem]{Proposition}
\newtheorem{lemma}[theorem]{Lemma}
\newtheorem{corollary}[theorem]{Corollary}

\theoremstyle{definition}
\newtheorem{definition}[theorem]{Definition}
\newtheorem{hypothesis}[theorem]{Hypothesis}

\theoremstyle{remark}
\newtheorem{remark}[theorem]{Remark}

\numberwithin{equation}{section}

\title[Minimal ancient solutions]{Minimal Ancient Solutions for the Navier--Stokes Equations in Self-similar Variables}

\author{Guillem Duran Ballester}
\date{}

\subjclass[2020]{35Q30, 35B44, 35B40, 76D05}
\keywords{Navier--Stokes equations, ancient solutions, Type I singularities, self-similar variables, concentration compactness}

\begin{document}

\begin{abstract}
We prove a minimal-element reduction for bounded ancient solutions of the
three-dimensional Navier--Stokes equations in backward self-similar variables.
The result is formulated with explicit compactness and closedness hypotheses:
if a nonempty time-translation invariant collection of bounded ancient
solutions has a fixed local \(L^3\) nonvanishing normalization and is closed
under locally smooth limits, then it contains a nonzero solution of minimal
\(L^\infty\) size.  Consequently, a Liouville theorem for such minimal
solutions implies the corresponding Liouville theorem for the whole collection.
We also record a concrete uniformly \(L^3\)-tight setting in which the required
closedness is verified directly.
\end{abstract}

\maketitle

\section{Introduction}

We consider the incompressible Navier--Stokes equations in backward
self-similar variables,
\begin{equation}\label{eq:centered}
   \partial_\tau V+(V\cdot\nabla)V+\nabla P
   =
   \Delta V-\frac12(V+y\cdot\nabla V),
   \qquad \operatorname{div}V=0,
\end{equation}
on \(\mathbb R^3\times\mathbb R\).  The variables are related to the usual
ancient Navier--Stokes variables by
\[
   V(y,\tau)=\sqrt{-t}\,U(y\sqrt{-t},t),\qquad \tau=-\log(-t),
   \qquad t<0.
\]
The drift in \eqref{eq:centered} fixes the self-similar center and scale.  The
remaining symmetry relevant to a Type I blow-up limit is translation of the
logarithmic time variable \(\tau\).

This paper isolates the compactness step which is needed before a rigidity
argument can be applied.  Suppose a family of bounded ancient solutions is
stable under time translation and locally smooth limits, and suppose every
solution in the family carries a definite amount of \(L^3\) mass in a fixed
ball at some time.  Among all such solutions, one can minimize the usual
\(L^\infty(\mathbb R^3\times\mathbb R)\) norm.  The infimum is positive and is
attained by a nonzero ancient solution.  The proof is short but important:
one first recenters the nonvanishing time, then uses local compactness and
lower semicontinuity of the \(L^\infty\) norm.

The main theorem is deliberately a reduction theorem.  It does not assert that
the resulting minimal solution is zero.  Rather, it states the precise
ordinary PDE conclusion which a subsequent rigidity argument must rule out.
In this form the statement uses only the standard compactness theory for
bounded ancient Navier--Stokes solutions.

\subsection{What must be checked}

The minimization argument uses four facts and no hidden compactness
assumption.  The class must be uniformly bounded in \(L^\infty\), otherwise
there is no compactness.  It must contain a fixed local nonvanishing
normalization, otherwise the minimizing sequence could converge to zero.  It
must be invariant under time translation, because the nonvanishing time of a
given solution need not be \(\tau=0\).  Finally, it must be closed under
locally smooth limits which retain the normalization.  These are exactly the
four clauses in \cref{hyp:closed-normalized}.

\section{Bounded ancient solutions}

After applying the Leray projection \(\mathbb P\), equation \eqref{eq:centered}
takes the form
\begin{equation}\label{eq:projected}
   \partial_\tau V
   =
   \Delta V-\frac12y\cdot\nabla V-\frac12V
   -\mathbb P\operatorname{div}(V\otimes V).
\end{equation}
Here \(\mathbb P=I-\nabla\Delta^{-1}\operatorname{div}\) is the usual
whole-space Helmholtz--Leray projection.  Thus \eqref{eq:projected} is
understood distributionally: for divergence-free test fields the pressure term
in \eqref{eq:centered} drops out, and the pressure is recovered, up to a
function of time, by the whole-space Riesz transform formula
\[
   -\Delta P=\partial_i\partial_j(V_iV_j).
\]
No local pressure gauge is used below; the mild formulation fixes this
whole-space Leray gauge.
Let \(S(a)\), \(a>0\), denote the semigroup generated by
\(\Delta-\frac12y\cdot\nabla-\frac12\).  It is given by
\begin{equation}\label{eq:OU-semigroup}
   (S(a)f)(y)
   =
   e^{-a/2}
   \int_{\mathbb R^3}
   \frac{\exp\!\left(-\frac{|z-e^{-a/2}y|^2}{4(1-e^{-a})}\right)}
        {[4\pi(1-e^{-a})]^{3/2}}
   f(z)\,dz .
\end{equation}
The kernel in \eqref{eq:OU-semigroup} has mass \(e^{-a/2}\).  Consequently
\[
   \|S(a)f\|_{L^\infty}\le e^{-a/2}\|f\|_{L^\infty},
   \qquad a>0.
\]
We shall also use the standard Oseen-kernel estimate, in self-similar
variables,
\begin{equation}\label{eq:stokes-adjoint-estimate}
   \|(S(a)\mathbb P\operatorname{div})^*\phi\|_{L^1(\mathbb R^3)}
   \le
   C e^{-a/2}(1-e^{-a})^{-1/2}\|\phi\|_{L^1(\mathbb R^3)}
\end{equation}
for \(\phi\in C_c^\infty(\mathbb R^3)\).  The singular factor is integrable at
\(a=0\), and for fixed \(a>0\) the adjoint kernel is smooth and rapidly
decaying when tested against a compactly supported smooth function.

\begin{definition}[Bounded mild ancient solution]\label{def:mild}
A vector field \(V:\mathbb R^3\times\mathbb R\to\mathbb R^3\) is a bounded mild
ancient solution of \eqref{eq:centered} if
\[
   V\in L^\infty(\mathbb R^3\times\mathbb R),\qquad \operatorname{div}V=0
\]
in distributions, and if for every \(\sigma<\tau\)
\begin{equation}\label{eq:mild}
   V(\tau)
   =
   S(\tau-\sigma)V(\sigma)
   -
   \int_\sigma^\tau
   S(\tau-s)\mathbb P\operatorname{div}(V\otimes V)(s)\,ds
\end{equation}
as an identity in the weak-* topology of \(L^\infty(\mathbb R^3)\).
\end{definition}

\begin{remark}
The mild formulation fixes the whole-space Leray pressure, modulo functions of
time.  This excludes the parasitic pressure terms which may otherwise appear
in local ancient solutions on the whole space.
\end{remark}

\begin{lemma}[Local smoothing]\label{lem:smoothing}
Let \(V\) be a bounded mild solution of \eqref{eq:centered} on
\(\mathbb R^3\times(\tau_0,\tau_1)\).  For every compact cylinder
\[
   K\Subset \mathbb R^3\times(\tau_0,\tau_1)
\]
and every integer \(m\ge0\),
\[
   \|V\|_{C^m(K)}
   \le
   C\bigl(m,K,\tau_0,\tau_1,\|V\|_{L^\infty(\mathbb R^3\times(\tau_0,\tau_1))}\bigr).
\]
The same estimate holds uniformly for any family with a common
\(L^\infty\) bound.
\end{lemma}

\begin{proof}
Fix compact intervals \(J_0\Subset J_1\Subset(\tau_0,\tau_1)\), with the time
projection of \(K\) contained in the interior of \(J_0\), and set
\(t=-e^{-\tau}\).  Since the map \(\tau\mapsto t\) is a smooth
diffeomorphism from \(J_1\) onto a compact interval
\(I_1\Subset(-\infty,0)\), define
\[
   U(x,t)=(-t)^{-1/2}V(x/\sqrt{-t},-\log(-t)).
\]
The change of variables is the standard backward self-similar transform.  A
direct chain-rule computation, first for smooth solutions and then by
approximation in the mild formulation, shows that \(U\) solves the usual
Navier--Stokes system
\[
   \partial_tU+(U\cdot\nabla)U+\nabla p=\Delta U,\qquad
   \operatorname{div}U=0
\]
on \(\mathbb R^3\times I_1\).  Moreover
\[
   \|U\|_{L^\infty(\mathbb R^3\times I_1)}
   \le
   C(J_1)\|V\|_{L^\infty(\mathbb R^3\times(\tau_0,\tau_1))},
\]
because \((-t)^{-1/2}\) is bounded above and below on \(I_1\).

Interior smoothing for bounded mild Navier--Stokes solutions, in the form of
Kato \cite{Kato1984} and Giga--Miyakawa \cite{GigaMiyakawa1985}, gives bounds
for all space-time derivatives of \(U\) on compact subcylinders of
\(\mathbb R^3\times I_0\), where \(I_0=t(J_0)\).  The constants depend only on
the derivative order, the compact cylinder, the distance from \(I_0\) to the
endpoints of \(I_1\), and the displayed \(L^\infty\) bound.  The map
\((y,\tau)\mapsto (y\sqrt{-t},t)\) and all of its derivatives are bounded on
the image of \(K\), so applying the chain rule back to \((y,\tau)\) gives the
stated estimate.  The same argument is uniform for families with a common
\(L^\infty\) bound.
\end{proof}

\begin{proposition}[Local compactness]\label{prop:local-compactness}
Let \(V_n\) be bounded mild ancient solutions of \eqref{eq:centered} satisfying
\[
   \sup_n\|V_n\|_{L^\infty(\mathbb R^3\times\mathbb R)}<\infty .
\]
Then a subsequence converges in
\[
   C^m_{\mathrm{loc}}(\mathbb R^3\times\mathbb R)
\]
for every \(m\ge0\) to a bounded mild ancient solution \(V\) of
\eqref{eq:centered}.
\end{proposition}

\begin{proof}
\Cref{lem:smoothing} gives uniform \(C^m\) bounds on every compact cylinder,
for every \(m\).  Choose an exhaustion \(Q_j=B_j\times[-j,j]\) of
\(\mathbb R^3\times\mathbb R\).  For each fixed \(j\) and \(m\),
Arzela--Ascoli gives precompactness in \(C^{m-1}(Q_j)\), and the uniform
\(C^m\) bounds identify the limit derivatives.  A diagonal extraction over the
countable pairs \((j,m)\) gives a subsequence, not relabeled, converging in
\(C^m(Q_j)\) for every \(j\) and every \(m\).  The limit \(V\) is smooth and
satisfies
\[
   \|V\|_{L^\infty(\mathbb R^3\times\mathbb R)}
   \le
   \sup_n\|V_n\|_{L^\infty(\mathbb R^3\times\mathbb R)}.
\]
The divergence-free condition passes to the limit locally in distributions.
The equation \eqref{eq:centered}, tested against compactly supported
space-time test functions, also passes to the limit because the convergence is
local \(C^1\) and the nonlinear products \(V_n\otimes V_n\) converge locally
uniformly to \(V\otimes V\).  Thus \(V\) is a smooth distributional solution of
\eqref{eq:centered}.

It remains to record that the mild formulation is closed under this
convergence.  Fix \(\sigma<\tau\).  We first note the elementary consequence of
local uniform convergence plus a common \(L^\infty\) bound: if
\(\psi\in L^1(\mathbb R^3)\), then
\[
   \int_{\mathbb R^3} V_n(y,t)\psi(y)\,dy
   \to
   \int_{\mathbb R^3} V(y,t)\psi(y)\,dy,
   \qquad t\in\{\sigma,\tau\}.
\]
Indeed, approximate \(\psi\) in \(L^1\) by a compactly supported function and
use the common \(L^\infty\) bound to control the tail uniformly in \(n\).  The
same argument applies to \(V_n\otimes V_n\), because
\(|V_n\otimes V_n|\le C\) uniformly and \(V_n\otimes V_n\to V\otimes V\)
locally uniformly.

Now test \eqref{eq:mild} for \(V_n\) against
\(\phi\in C_c^\infty(\mathbb R^3)\).  The adjoint of \(S(\tau-\sigma)\) maps
\(\phi\) to an \(L^1\) function, so the linear term passes to the limit by the
previous paragraph.  For the Duhamel term, duality and
\eqref{eq:stokes-adjoint-estimate} give
\[
   \|(S(\tau-s)\mathbb P\operatorname{div})^*\phi\|_{L^1}
   \le
   C_\phi e^{-(\tau-s)/2}(1-e^{-(\tau-s)})^{-1/2}.
\]
The right-hand side is integrable in \(s\in(\sigma,\tau)\).  For each fixed
\(s<\tau\), the preceding \(L^1\)-duality argument gives convergence of the
integrand with \(V_n\otimes V_n\) replaced by \(V\otimes V\), while the same
estimate supplies an integrable dominating function.  Dominated convergence
therefore passes the Duhamel term to the limit.  Hence \eqref{eq:mild} holds
for \(V\) against all compactly supported smooth test functions.

Both sides of \eqref{eq:mild} are bounded elements of \(L^\infty(\mathbb R^3)\):
the linear term is bounded by the semigroup estimate, and the Duhamel term is
bounded weak-* by the preceding \(L^1\)-adjoint estimate and the common
\(L^\infty\) bound.  Since \(C_c^\infty(\mathbb R^3)\) is dense in
\(L^1(\mathbb R^3)\), equality against compactly supported smooth tests
extends to equality against every \(L^1\) test function.  This is precisely the
weak-* \(L^\infty\) identity \eqref{eq:mild}.  Thus the limit is a bounded mild
ancient solution.
\end{proof}

\section{Closed normalized collections}

The only symmetry used below is translation in the time variable.  For a
solution \(V\) and \(s\in\mathbb R\), write
\[
   V^{s}(y,\tau)=V(y,\tau+s).
\]

\begin{hypothesis}[Closed normalized collection]\label{hyp:closed-normalized}
Let \(A\) be a collection of bounded mild ancient solutions of
\eqref{eq:centered}.  Assume the following.

\begin{enumerate}
\item There is \(M<\infty\) such that
\[
   \|V\|_{L^\infty(\mathbb R^3\times\mathbb R)}\le M
   \qquad\hbox{for every }V\in A .
\]

\item There are \(R_0>0\) and \(\eta_0>0\) such that
\begin{equation}\label{eq:local-nonzero}
   \sup_{\tau\in\mathbb R}\int_{B_{R_0}}|V(y,\tau)|^3\,dy\ge\eta_0
   \qquad\hbox{for every }V\in A .
\end{equation}

\item If \(V\in A\) and \(s\in\mathbb R\), then \(V^s\in A\).

\item If \(V_n\in A\), \(V_n\to V\) locally smoothly on
\(\mathbb R^3\times\mathbb R\), and \(V\) satisfies \eqref{eq:local-nonzero},
then \(V\in A\).
\end{enumerate}
\end{hypothesis}

\begin{remark}
The constants \(R_0\) and \(\eta_0\) are part of the normalization.  In a
Type I blow-up construction they come from the scale selection which prevents
the ancient limit from vanishing.
\end{remark}

\begin{lemma}[Time centering]\label{lem:time-centering}
Let \(A\) satisfy \cref{hyp:closed-normalized}.  If \(V\in A\), then for every
\(\delta>0\) there is \(s\in\mathbb R\) such that \(V^s\in A\) and
\[
   \int_{B_{R_0}}|V^s(y,0)|^3\,dy\ge\eta_0-\delta .
\]
\end{lemma}

\begin{proof}
The quantity
\[
   \sup_{\tau\in\mathbb R}\int_{B_{R_0}}|V(y,\tau)|^3\,dy
\]
need not be attained.  The definition of the supremum nevertheless implies
that, for every \(\delta>0\), there is \(s\in\mathbb R\) such that
\[
   \int_{B_{R_0}}|V(y,s)|^3\,dy\ge\eta_0-\delta .
\]
Then \(V^s\in A\) by time-translation invariance, and the displayed inequality
is exactly the asserted one at \(\tau=0\).
\end{proof}

\begin{lemma}[Stability of the normalization]\label{lem:nonzero-stability}
Assume \(V_n\to V\) locally uniformly on \(\mathbb R^3\times\mathbb R\).  If
\[
   \int_{B_{R_0}}|V_n(y,0)|^3\,dy\ge \eta_0-\delta_n,
   \qquad \delta_n\downarrow0,
\]
then
\[
   \int_{B_{R_0}}|V(y,0)|^3\,dy\ge \eta_0 .
\]
\end{lemma}

\begin{proof}
Restrict the convergence to the compact set \(B_{R_0}\times\{0\}\).  Uniform
convergence there implies
\[
   \|V_n(\cdot,0)-V(\cdot,0)\|_{L^3(B_{R_0})}
   \le
   |B_{R_0}|^{1/3}
   \|V_n(\cdot,0)-V(\cdot,0)\|_{L^\infty(B_{R_0})}
   \to0 .
\]
By continuity of the map \(f\mapsto\int_{B_{R_0}}|f|^3\) on
\(L^3(B_{R_0})\),
\[
   \int_{B_{R_0}}|V(y,0)|^3\,dy
   =
   \lim_{n\to\infty}\int_{B_{R_0}}|V_n(y,0)|^3\,dy
   \ge\eta_0 .
\]
\end{proof}

\begin{lemma}[Lower semicontinuity of size]\label{lem:lsc}
Assume \(V_n\to V\) locally uniformly on \(\mathbb R^3\times\mathbb R\).  Then
\[
   \|V\|_{L^\infty(\mathbb R^3\times\mathbb R)}
   \le
   \liminf_{n\to\infty}
   \|V_n\|_{L^\infty(\mathbb R^3\times\mathbb R)} .
\]
\end{lemma}

\begin{proof}
If the liminf is infinite, there is nothing to prove.
Choose a subsequence \(n_k\) such that
\[
   \|V_{n_k}\|_{L^\infty(\mathbb R^3\times\mathbb R)}
   \to
   L:=
   \liminf_{n\to\infty}
   \|V_n\|_{L^\infty(\mathbb R^3\times\mathbb R)} .
\]
For every point \((y,\tau)\), local uniform convergence gives pointwise
convergence along this subsequence, hence
\[
   |V(y,\tau)|
   =
   \lim_{k\to\infty}|V_{n_k}(y,\tau)|
   \le
   \liminf_{k\to\infty}
   \|V_{n_k}\|_{L^\infty(\mathbb R^3\times\mathbb R)}
   =L .
\]
Taking the supremum over \((y,\tau)\) gives the claim.
\end{proof}

\section{The minimal solution}

\begin{theorem}[Existence of a minimal ancient solution]\label{thm:minimal-existence}
Let \(A\) be nonempty and satisfy \cref{hyp:closed-normalized}.  Set
\[
   m=\inf\left\{
      \|V\|_{L^\infty(\mathbb R^3\times\mathbb R)}:\ V\in A
   \right\}.
\]
Then
\[
   0<m\le M,
\]
and there exists \(V_*\in A\) such that
\[
   \|V_*\|_{L^\infty(\mathbb R^3\times\mathbb R)}=m .
\]
In particular \(V_*\not\equiv0\).
\end{theorem}

\begin{proof}
The upper bound \(m\le M\) follows from \cref{hyp:closed-normalized}.  The
local nonvanishing gives a positive lower bound.  Indeed, for every \(V\in A\),
\[
   \eta_0
   \le
   \sup_{\tau\in\mathbb R}\int_{B_{R_0}}|V(y,\tau)|^3\,dy
   \le
   |B_{R_0}|\,
   \|V\|_{L^\infty(\mathbb R^3\times\mathbb R)}^3 .
\]
Thus
\[
   m\ge \left(\frac{\eta_0}{|B_{R_0}|}\right)^{1/3}>0 .
\]

Choose \(V_n\in A\) such that
\[
   m
   \le
   \|V_n\|_{L^\infty(\mathbb R^3\times\mathbb R)}
   \le
   m+\frac1n .
\]
For each \(n\), apply \cref{lem:time-centering} with \(\delta=1/n\).  Replacing
\(V_n\) by the corresponding time translate keeps \(V_n\) inside \(A\) and
does not change its \(L^\infty\) norm, since
\[
   \sup_{(y,\tau)\in\mathbb R^3\times\mathbb R}|V_n(y,\tau+s)|
   =
   \sup_{(y,\tau)\in\mathbb R^3\times\mathbb R}|V_n(y,\tau)|.
\]
Thus we may assume that
\[
   \int_{B_{R_0}}|V_n(y,0)|^3\,dy\ge\eta_0-\frac1n .
\]
\Cref{prop:local-compactness} applies because all \(V_n\) obey the common
\(L^\infty\) bound \(M\).  Hence a subsequence, not relabeled, converges
locally smoothly to a bounded mild ancient solution \(V_*\).
By \cref{lem:nonzero-stability},
\[
   \int_{B_{R_0}}|V_*(y,0)|^3\,dy\ge\eta_0 .
\]
In particular \(V_*\) satisfies \eqref{eq:local-nonzero}.  The closedness
assumption in \cref{hyp:closed-normalized} therefore gives \(V_*\in A\).
Finally, by \cref{lem:lsc},
\[
   \|V_*\|_{L^\infty(\mathbb R^3\times\mathbb R)}
   \le
   \liminf_{n\to\infty}
   \|V_n\|_{L^\infty(\mathbb R^3\times\mathbb R)}
   =
   m .
\]
Since \(m\) is the infimum over all elements of \(A\) and \(V_*\in A\), one
also has \(m\le\|V_*\|_{L^\infty}\).  The two inequalities give equality.
The lower bound \(m>0\) shows that the minimizer cannot be the zero solution.
\end{proof}

\begin{corollary}[Reduction to minimal solutions]\label{cor:minimal-reduction}
Let \(A\) satisfy \cref{hyp:closed-normalized}.  If an independent argument
rules out every nonzero element of \(A\) which attains the minimum in
\cref{thm:minimal-existence}, then \(A\) is empty.
\end{corollary}

\begin{proof}
This is the contrapositive of \cref{thm:minimal-existence}.
\end{proof}

\section{A uniformly tight setting}

We record a standard way to verify the closedness hypothesis.  Fix positive
constants \(M\), \(L\), \(R_0\), and \(\eta_0\), and let
\(\omega:[1,\infty)\to[0,\infty)\) be a function with
\[
   \lim_{R\to\infty}\omega(R)=0 .
\]
Let \(A\) be the collection of all bounded mild ancient solutions of
\eqref{eq:centered} satisfying
\begin{equation}\label{eq:tight-class-bounds}
   \|V\|_{L^\infty(\mathbb R^3\times\mathbb R)}\le M,
   \qquad
   \sup_{\tau\in\mathbb R}\|V(\tau)\|_{L^3(\mathbb R^3)}\le L,
\end{equation}
\begin{equation}\label{eq:tight-class-tail}
   \sup_{\tau\in\mathbb R}
   \int_{|y|>R}|V(y,\tau)|^3\,dy\le\omega(R)
   \qquad\hbox{for every }R\ge1,
\end{equation}
and the local nonvanishing condition \eqref{eq:local-nonzero}.

\begin{proposition}[Closedness under uniform tightness]\label{prop:tight-closed}
The collection \(A\) defined by \eqref{eq:tight-class-bounds},
\eqref{eq:tight-class-tail}, and \eqref{eq:local-nonzero} satisfies
\cref{hyp:closed-normalized}.
\end{proposition}

\begin{proof}
The \(L^\infty\) bound and the local nonvanishing condition are part of the
definition of \(A\).  Time translation preserves \eqref{eq:centered},
\eqref{eq:tight-class-bounds}, \eqref{eq:tight-class-tail}, and
\eqref{eq:local-nonzero}.

It remains to verify closedness.  Let \(V_n\in A\) and suppose
\[
   V_n\to V
   \quad\hbox{locally smoothly on }\mathbb R^3\times\mathbb R,
\]
where \(V\) satisfies \eqref{eq:local-nonzero}.  By
\cref{prop:local-compactness}, any locally smooth limit of a subsequence of
\((V_n)\) is a bounded mild ancient solution.  The given locally smooth limit
is unique on every compact cylinder, so \(V\) itself is that bounded mild
ancient solution.  In particular the equation and the Leray mild formulation
hold for \(V\).
\Cref{lem:lsc} gives the \(L^\infty\) bound in
\eqref{eq:tight-class-bounds}.

For the \(L^3\) bound, fix \(R<\infty\) and \(\tau\in\mathbb R\).  Local smooth
convergence gives
\[
   \int_{B_R}|V(y,\tau)|^3\,dy
   =
   \lim_{n\to\infty}\int_{B_R}|V_n(y,\tau)|^3\,dy
   \le L^3 .
\]
The integrals on the left are monotone in \(R\).  By the monotone convergence
theorem,
\[
   \|V(\tau)\|_{L^3(\mathbb R^3)}^3
   =
   \lim_{R\to\infty}\int_{B_R}|V(y,\tau)|^3\,dy
   \le L^3 .
\]
This holds for every fixed \(\tau\), and therefore
\[
   \sup_{\tau\in\mathbb R}\|V(\tau)\|_{L^3(\mathbb R^3)}\le L .
\]
For the tail estimate, fix \(1\le R<R'<\infty\).  Then for every \(\tau\),
\[
   \int_{R<|y|<R'}|V(y,\tau)|^3\,dy
   =
   \lim_{n\to\infty}\int_{R<|y|<R'}|V_n(y,\tau)|^3\,dy
   \le \omega(R).
\]
The annuli \(\{R<|y|<R'\}\) increase to \(\{|y|>R\}\) as
\(R'\to\infty\).  Applying monotone convergence again yields
\[
   \int_{|y|>R}|V(y,\tau)|^3\,dy\le\omega(R).
\]
Taking the supremum in \(\tau\) proves \eqref{eq:tight-class-tail} for \(V\).
The assumed local nonvanishing \eqref{eq:local-nonzero} supplies the final
defining condition of \(A\).  Thus \(V\in A\), which is exactly the closedness
clause in \cref{hyp:closed-normalized}.
\end{proof}

\begin{corollary}[Minimal solution in the uniformly tight setting]\label{cor:tight-minimal}
If the collection \(A\) in \cref{prop:tight-closed} is nonempty, then \(A\)
contains a nonzero solution of minimal
\[
   L^\infty(\mathbb R^3\times\mathbb R)
\]
norm.
\end{corollary}

\begin{proof}
This follows from \cref{prop:tight-closed} and
\cref{thm:minimal-existence}.
\end{proof}

\begin{corollary}[Interface with stationary rigidity]\label{cor:stationary-interface}
Let \(A\) be the uniformly tight collection in \cref{prop:tight-closed}.  If
every minimal element of \(A\) is stationary in \(\tau\), then \(A\) is empty.
\end{corollary}

\begin{proof}
If \(A\) were nonempty, \cref{cor:tight-minimal} would give a nonzero minimal
element \(V_*\).  By assumption \(V_*(y,\tau)=W(y)\).  The bounds in
\eqref{eq:tight-class-bounds} imply \(W\in L^3(\mathbb R^3)\), and the
\(L^\infty\) bound and local smoothing imply that \(W\) is a smooth bounded
solution.  Substituting \(\partial_\tau V_*=0\) into \eqref{eq:centered}, or
equivalently testing the time-independent distributional equation against
compactly supported fields, gives the stationary
equation is
\[
   (W\cdot\nabla)W+\nabla P
   =
   \Delta W-\frac12(W+y\cdot\nabla W),
   \qquad \operatorname{div}W=0 .
\]
The theorem of Necas--Ruzicka--Sverak \cite{NRS1996} then gives \(W\equiv0\),
contradicting \eqref{eq:local-nonzero}.
\end{proof}

\section{Conclusion}

The result of this paper is the following precise reduction.  Under the
closedness, time-translation invariance, boundedness, and fixed local
nonvanishing hypotheses stated above, any nonempty collection of bounded
ancient solutions contains a nonzero solution of smallest
\(L^\infty(\mathbb R^3\times\mathbb R)\) norm.  Therefore a full Liouville
theorem in such a collection is reduced to the exclusion of these minimal
solutions.  In the uniformly \(L^3\)-tight setting, the hypotheses are closed
under locally smooth limits, and stationary minimal solutions are already
excluded by the stationary self-similar theorem.

\begin{thebibliography}{99}

\bibitem{GigaMiyakawa1985}
Y.~Giga and T.~Miyakawa,
\emph{Solutions in \(L^r\) of the Navier--Stokes initial value problem},
Arch. Rational Mech. Anal. \textbf{89} (1985), 267--281.

\bibitem{Kato1984}
T.~Kato,
\emph{Strong \(L^p\)-solutions of the Navier--Stokes equation in
\(\mathbb R^m\), with applications to weak solutions},
Math. Z. \textbf{187} (1984), 471--480.

\bibitem{NRS1996}
J.~Necas, M.~Ruzicka, and V.~Sverak,
\emph{On Leray's self-similar solutions of the Navier--Stokes equations},
Acta Math. \textbf{176} (1996), 283--294.

\end{thebibliography}

\end{document}
